% ============================================================================
% templates/report/cornell/cornell.tex — plantilla: hoja Cornell de apuntes
% ----------------------------------------------------------------------------
% Página 1: la hoja en blanco para imprimir (preguntas, notas, resumen).
% Página 2: un ejemplo lleno (Econometría I, diseño muestral de la ENAHO).
% Página 3: por qué y cómo se usa, con sus referencias.
% La misma plantilla para la web: cornell.md. SIN diseño: lo pone academic-report.
% Uso: scripts/build.sh cornell.tex   (LuaLaTeX)
% ============================================================================
\documentclass{academic-report}
\tipodocumento{Hoja Cornell}
\universidad{Universidad Nacional de San Cristóbal de Huamanga}
\escuela{Escuela Profesional de Economía}
\curso{Nombre del curso}\docente{Edison Achalma B.Sc. Econ.}\periodo{2026-I}
\titulodocumento{Apuntes de la sesión}

\begin{document}
\encabezado

% ------------------------------------------------------- 1. la hoja vacía ---
\noindent
\begin{tabularx}{\linewidth}{|X|X|X|}
\hline
\textbf{Curso:} & \textbf{Sesión:} & \textbf{Fecha:} \\[0.5em]
\hline
\end{tabularx}

\vspace{0.6em}\noindent
\begin{tabularx}{\linewidth}{|p{0.27\linewidth}|X|}
\hline
\textbf{Preguntas} {\small(después de clase)} & \textbf{Notas} {\small(durante la clase)} \\
\hline
\rule{0pt}{12.6cm} & \\
\hline
\multicolumn{2}{|l|}{\textbf{Resumen} {\small(después de clase, con tus palabras)}} \\
\multicolumn{2}{|l|}{\rule{0pt}{3.4cm}} \\
\hline
\end{tabularx}

% ------------------------------------------------------ 2. ejemplo lleno ---
\newpage
\section*{Ejemplo lleno}
\noindent
\begin{tabularx}{\linewidth}{|X|X|X|}
\hline
\textbf{Curso:} Econometría I & \textbf{Sesión:} Diseño muestral de la ENAHO & \textbf{Fecha:} \\
\hline
\end{tabularx}

\vspace{0.6em}\noindent
{\small
\begin{tabularx}{\linewidth}{|>{\raggedright\arraybackslash}p{0.27\linewidth}|>{\raggedright\arraybackslash}X|}
\hline
\textbf{Preguntas} & \textbf{Notas} \\
\hline
¿Por qué la ENAHO no es una muestra aleatoria simple (MAS)?
& Muestra probabilística, de áreas, estratificada, bietápica e independiente en cada
  departamento (INEI, 2026, p.~404). Cada hogar tiene probabilidad conocida de entrar,
  pero no la misma. \\[0.4em]
¿Qué es un estrato y qué gana el diseño con él?
& \textbf{Estrato}: grupo que no se mezcla con otro; se sortea dentro de cada uno.
  Asegura que todos los dominios estén en la muestra; suele bajar la varianza. \\[0.4em]
¿Qué se sortea en cada etapa? ¿Qué cuesta sortear conglomerados?
& \textbf{Conglomerado}: grupo de viviendas vecinas. 1.ª etapa: conglomerados;
  2.ª etapa: viviendas dentro de cada uno. Abarata el campo, pero los vecinos se
  parecen (correlación dentro del conglomerado): cada hogar aporta menos información. \\[0.4em]
¿Qué representa el factor de expansión? ¿Por qué sale 17,2\,\% y no 20,2\,\%?
& \textbf{Factor de expansión} (\texttt{factor07}): cuántos hogares del país representa
  cada hogar de la muestra; parte de la inversa de la probabilidad de selección.
  Pobreza de personas: peso \texttt{factor07} $\times$ \texttt{mieperho}.
  ENAHO 2019 sin pesos y sobre hogares: 17,2\,\%. Con pesos y sobre personas: 20,2\,\%,
  la cifra oficial (INEI, 2026, p.~72). \\[0.4em]
¿Qué mide el efecto de diseño? ¿Qué pasa si lo ignoro?
& \textbf{Efecto de diseño}: $\mathit{deff} = \operatorname{Var}_{\text{diseño}} / \operatorname{Var}_{\text{MAS}}$.
  Tasa de 2019: EE como MAS 0,22~pp; EE con el diseño 0,44~pp (el doble);
  $\mathit{deff} \approx 4{,}2$ con los EE sin redondear.
  Ignorarlo da intervalos la mitad de anchos: falsa precisión, diferencias
  “significativas” que no lo son. Muestra efectiva $\approx n/\mathit{deff}$. \\[0.4em]
¿Cómo se declara el diseño antes de estimar?
& Stata: \texttt{svyset conglome [pw=factor07], strata(estrato)} y luego
  \texttt{svy: mean pobre} (\texttt{pobre} = 1 si \texttt{pobreza} $\leq 2$).
  Sin \texttt{svyset}, el EE sale como MAS. \\[0.4em]
\hline
\multicolumn{2}{|p{\dimexpr\linewidth-2\tabcolsep-2\arrayrulewidth\relax}|}{\textbf{Resumen.}
  La ENAHO no es una MAS: estratifica, sortea conglomerados y luego viviendas, y pondera.
  Sin el factor de expansión y contando hogares la tasa sale 17,2\,\% y no el 20,2\,\% oficial;
  sin el diseño, el error estándar sale la mitad del verdadero ($\mathit{deff} \approx 4{,}2$).
  Toda estimación con la ENAHO se declara con estratos, conglomerados y pesos.} \\
\hline
\end{tabularx}}
\notatabla{Cifras de pobreza: INEI (2026, p.~72). Tasa sin pesos, errores estándar y efecto de
diseño: cálculo propio con la ENAHO 2019 (pp = puntos porcentuales; EE = error estándar).}

% ------------------------------------------------ 3. por qué y cómo se usa ---
\newpage
\section*{Por qué}
Los apuntes rinden sobre todo cuando se repasan después, no solo cuando se toman: el
beneficio de tomar notas está en buena parte en su revisión posterior (Kiewra, 1985). Por eso
la hoja deja una columna para preguntas: es el insumo del repaso, y lo que no sale al
responderlas vuelve a estudiarse.

\begin{nota}
La columna de preguntas no se llena en clase: se escribe después, a partir de las notas, y es
la que se usa para repasar (taparse las notas y responder). Esta hoja no supone que escribir a
mano sea mejor que en teclado; funciona igual en papel (este PDF) y en la web (\texttt{cornell.md}).
\end{nota}

\section*{Cómo se usa}
La hoja original reparte la página en una columna de preguntas de 2½ pulgadas, una de notas de
6 y una franja de resumen de 2 (Cornell University, Learning Strategies Center, 2013, p.~1); esta
la conserva en proporción sobre A4. Sus cinco pasos:
\begin{enumerate}
  \item \textbf{Registrar.} En clase, en la columna de notas, con frases cortas.
  \item \textbf{Preguntar.} Lo antes posible tras la clase, escribe a la izquierda las
        preguntas que las notas responden.
  \item \textbf{Recitar.} Tapa las notas y responde cada pregunta en voz alta, con tus palabras.
  \item \textbf{Reflexionar.} Pregúntate qué significan esos hechos, en qué principio se
        apoyan y cómo se aplican.
  \item \textbf{Repasar.} Dedica al menos diez minutos a la semana a repasar las hojas
        anteriores (Cornell University, Learning Strategies Center, 2013, p.~1).
\end{enumerate}
Al terminar, el resumen de abajo dice en dos o tres oraciones lo que la página enseña.

\noindent\textbf{Fuente:} Cornell University, Learning Strategies Center (2013, Calibre 11126);
Kiewra (1985).

\section*{Referencias}
{\setlength{\parindent}{0pt}
\hangindent=1.27cm\hangafter=1
Cornell University, Learning Strategies Center. (2013). \textit{The Cornell note-taking system}.\par
\hangindent=1.27cm\hangafter=1
Instituto Nacional de Estadística e Informática. (2026). \textit{Perú: Evolución de la pobreza
monetaria 2016-2025. Informe técnico}. INEI.\par
\hangindent=1.27cm\hangafter=1
Kiewra, K. A. (1985). Investigating notetaking and review: A depth of processing alternative.
\textit{Educational Psychologist, 20}(1), 23–32. \url{https://doi.org/10.1207/s15326985ep2001_4}\par}

\end{document}
